% GENERATED FILE. Edit canonical lessons, then run npm run generate:books.
% canonical-source-hash: fnv1a64:77096018284c9663
% canonical-lessons: BN-C232-shosa, BN-C232-mashrum, BN-C232-bhutta, BN-C232-chhola, BN-C232-intarnet

\chapter{Cucumber, Mushroom, Maize, Chickpeas, Internet}
\label{ch:bn-cucumber-mushroom-maize-chickpeas-internet}
\hlchaptermodality{232}

\begin{chapteropening}
\textbf{By the end of this chapter:} \emph{I can say a cucumber, a mushroom, maize, chickpeas and the internet.}

Complete the last lesson of chapter 232: I can say a cucumber, a mushroom, maize, chickpeas and the internet.
\end{chapteropening}

\section[\texorpdfstring{\bn{শসা} (shôsā)}{শসা (shôsā)}]{\bn{শসা} (shôsā) — a cucumber}

\label{lesson:BN-C232-shosa}



\begin{quote}
Before the new one: say the Bengali for fine flour, then the Bengali for semolina.
\end{quote}

\subsection*{You'll want to know: \bn{শসা}}
\textbf{\bn{শসা}} — \emph{shôsā} — \textquotedblleft{}a cucumber\textquotedblright{}.

\begin{grammarlens}[title={What you've built}]
One more word for this chapter.
\end{grammarlens}

\subsection*{Your turn}
\begin{itemize}
\raggedright
  \item \emph{Say it:} \emph{shôsā}
  \item \emph{Say it:} \emph{shôsā}, once more
  \item \emph{Say it:} say \emph{shôsā}
  \item \emph{Recall:} say the Bengali for fine flour, then the Bengali for semolina, then say \emph{shôsā} again
\end{itemize}

\begin{tcolorbox}[breakable,colback=teal!4,colframe=teal!35!black,title={Before you move on}]
What does \bn{শসা} mean? (\textquotedblleft{}A cucumber\textquotedblright{}.) Say it once more.
\end{tcolorbox}
\section[\texorpdfstring{\bn{মাশরুম} (māshrum)}{মাশরুম (māshrum)}]{\bn{মাশরুম} (māshrum) — a mushroom}

\label{lesson:BN-C232-mashrum}



\begin{quote}
Before the new one: say the Bengali for semolina, then the Bengali for a cucumber.
\end{quote}

\subsection*{You'll want to know: \bn{মাশরুম}}
\textbf{\bn{মাশরুম}} — \emph{māshrum} — \textquotedblleft{}a mushroom\textquotedblright{}.

\begin{grammarlens}[title={What you've built}]
One more word for this chapter.
\end{grammarlens}

\subsection*{Your turn}
\begin{itemize}
\raggedright
  \item \emph{Say it:} \emph{māshrum}
  \item \emph{Say it:} \emph{māshrum}, once more
  \item \emph{Say it:} say \emph{māshrum}
  \item \emph{Recall:} say the Bengali for semolina, then the Bengali for a cucumber, then say \emph{māshrum} again
\end{itemize}

\begin{tcolorbox}[breakable,colback=teal!4,colframe=teal!35!black,title={Before you move on}]
What does \bn{মাশরুম} mean? (\textquotedblleft{}A mushroom\textquotedblright{}.) Say it once more.
\end{tcolorbox}
\section[\texorpdfstring{\bn{ভুট্টা} (bhuṭṭā)}{ভুট্টা (bhuṭṭā)}]{\bn{ভুট্টা} (bhuṭṭā) — maize, a corn cob}

\label{lesson:BN-C232-bhutta}



\begin{quote}
Before the new one: say the Bengali for a cucumber, then the Bengali for a mushroom.
\end{quote}

\subsection*{You'll want to know: \bn{ভুট্টা}}
\textbf{\bn{ভুট্টা}} — \emph{bhuṭṭā} — \textquotedblleft{}maize, a corn cob\textquotedblright{}.

\begin{grammarlens}[title={What you've built}]
One more word for this chapter.
\end{grammarlens}

\subsection*{Your turn}
\begin{itemize}
\raggedright
  \item \emph{Say it:} \emph{bhuṭṭā}
  \item \emph{Say it:} \emph{bhuṭṭā}, once more
  \item \emph{Say it:} say \emph{bhuṭṭā}
  \item \emph{Recall:} say the Bengali for a cucumber, then the Bengali for a mushroom, then say \emph{bhuṭṭā} again
\end{itemize}

\begin{tcolorbox}[breakable,colback=teal!4,colframe=teal!35!black,title={Before you move on}]
What does \bn{ভুট্টা} mean? (\textquotedblleft{}Maize, a corn cob\textquotedblright{}.) Say it once more.
\end{tcolorbox}
\section[\texorpdfstring{\bn{ছোলা} (chholā)}{ছোলা (chholā)}]{\bn{ছোলা} (chholā) — chickpeas}

\label{lesson:BN-C232-chhola}



\begin{quote}
Before the new one: say the Bengali for a mushroom, then the Bengali for maize.
\end{quote}

\subsection*{You'll want to know: \bn{ছোলা}}
\textbf{\bn{ছোলা}} — \emph{chholā} — \textquotedblleft{}chickpeas\textquotedblright{}.

\begin{grammarlens}[title={What you've built}]
One more word for this chapter.
\end{grammarlens}

\subsection*{Your turn}
\begin{itemize}
\raggedright
  \item \emph{Say it:} \emph{chholā}
  \item \emph{Say it:} \emph{chholā}, once more
  \item \emph{Say it:} say \emph{chholā}
  \item \emph{Recall:} say the Bengali for a mushroom, then the Bengali for maize, then say \emph{chholā} again
\end{itemize}

\begin{tcolorbox}[breakable,colback=teal!4,colframe=teal!35!black,title={Before you move on}]
What does \bn{ছোলা} mean? (\textquotedblleft{}Chickpeas\textquotedblright{}.) Say it once more.
\end{tcolorbox}
\section[\texorpdfstring{\bn{ইন্টারনেট} (inṭārneṭ)}{ইন্টারনেট (inṭārneṭ)}]{\bn{ইন্টারনেট} (inṭārneṭ) — the internet}

\label{lesson:BN-C232-intarnet}



\begin{quote}
Before the new one: say the Bengali for maize, then the Bengali for chickpeas.
\end{quote}

\subsection*{You'll want to know: \bn{ইন্টারনেট}}
\textbf{\bn{ইন্টারনেট}} — \emph{inṭārneṭ} — \textquotedblleft{}the internet\textquotedblright{}.

\begin{grammarlens}[title={What you've built}]
One more word for this chapter.
\end{grammarlens}

\subsection*{Your turn}
\begin{itemize}
\raggedright
  \item \emph{Say it:} \emph{inṭārneṭ}
  \item \emph{Say it:} \emph{inṭārneṭ}, once more
  \item \emph{Say it:} say \emph{inṭārneṭ}
  \item \emph{Recall:} say the Bengali for maize, then the Bengali for chickpeas, then say \emph{inṭārneṭ} again
\end{itemize}

\begin{tcolorbox}[breakable,colback=teal!4,colframe=teal!35!black,title={Before you move on}]
What does \bn{ইন্টারনেট} mean? (\textquotedblleft{}The internet\textquotedblright{}.) Say it once more.
\end{tcolorbox}
